\documentclass[10pt,a4paper]{amsart}
\usepackage[margin=19mm]{geometry}
\usepackage{amsmath,amssymb,mathtools}
\usepackage[scr=rsfs,cal=euler]{mathalfa}
\usepackage{xfrac}
\usepackage[unicode,hidelinks,bookmarks=false]{hyperref}
\newcommand{\ie}{i.e.\ }
\newcommand{\cf}{cf.\ }
\newcommand{\resp}{respectively\ }
\newcommand{\Subsection}{\subsection}
\providecommand{\nobreakdash}{\nobreak\hbox{-}}
\setlength{\emergencystretch}{2em}
\multlinegap0pt
\usepackage{bm}
\usepackage{bbm}
\usepackage{dsfont}
\usepackage{xspace}
\usepackage{cancel}
\usepackage{mathtools}
\usepackage{amsthm}
\makeatletter
\@ifundefined{lem}{\newtheorem{lem}{Lemma}}{}
\@ifundefined{prop}{\newtheorem{prop}{Proposition}}{}
\makeatother
\title[Achiral Lefschetz fibrations and the moduli space of curves]{Achiral Lefschetz fibrations and the moduli space of curves}
\author{Sardor Yakupov}
\date{Source: arXiv:2510.18593v1 (CC BY 4.0)}
\begin{document}
\maketitle

\noindent\emph{Source and licence.} Achiral Lefschetz fibrations and the moduli space of curves. Version arXiv:2510.18593v1 (CC BY 4.0), distributed under the Creative Commons Attribution 4.0 International licence, as recorded in the arXiv licence metadata for this exact version and shown on the abstract page https://arxiv.org/abs/2510.18593v1. Full rights notice: licenses/arxiv-2510.18593-CC-BY-4.0.txt.\par\smallskip

\noindent\textit{Editorial scope.} Counted unit: the Lem whose source label is \texttt{None} in the article (source0.tex lines 239--242, source label \texttt{None}), delivered together with the article's own complete original proof of it (source0.tex lines 243--312, one \texttt{proof} environment of 70 lines). No mathematics is written, repaired or re-derived: statement and proof are the article's own text line for line. Required context retained uncounted: maxprinc (lines 177--183). Every definition, display and label of those lines is retained unchanged. Omitted from this package and not redistributed: 1--176, 184--238, 313--453. Nothing inside a retained excerpt is omitted or altered: every retained body is delivered exactly as the article's lines are written, so no omission receipt of any kind accompanies this package. Citations: the retained excerpts carry no citation command at all, so no bibliography entry of the article is transcribed and no reference is redistributed. Reference adaptation: none was needed and none was made; every reference inside the delivered unit resolves inside it, so no unresolved reference or citation is produced. Printed slips kept literal: no printed word, symbol, hypothesis or number of the counted statement or of its proof was altered. Layout and class: the article's own class is \texttt{article} with packages including graphicx, amsmath, amssymb, amsthm, subfig, tikz, biblatex; the wrapper is the lane's pinned \texttt{amsart} editorial wrapper at 10pt on A4, which restates the theorem-like environments the retained text needs and adds the article's own macro definitions that the retained text needs (none), with extra wrapper packages (bm, bbm, dsfont, xspace, cancel, mathtools). The delivered numbering is the unit's own. Assets: no retained span contains a figure environment, a graphics-inclusion command, a PDF-page inclusion, a table or a TikZ picture; no image file is copied into this package, the package declares no asset dependency and no image file is redistributed with it. Licence: version arXiv:2510.18593v1 of this article is distributed under the Creative Commons Attribution 4.0 International licence, as recorded in the arXiv licence metadata for this exact version and shown on the abstract page https://arxiv.org/abs/2510.18593v1. A full rights notice is delivered at licenses/arxiv-2510.18593-CC-BY-4.0.txt. Characters: every retained excerpt is pure ASCII, so the delivered engine, XeLaTeX with its default Latin Modern text font, reports no missing character.
\begin{prop}[Maximum principle]\label{maxprinc}
    Let $(M, g_t)$ be a closed $n$-manifold equipped with a smooth family of Riemannian metrics. Consider $u:M \to \mathbb{R}$ a smooth function satisfying 
    $$\partial_tu - \Delta_{g_t}u \leq g_t(V(t), \nabla^{g_t}u)+F(u)$$
    for some vector field $V(t)$ and Lipschitz function $F$. Then for any constant $c$ such that $u_{t=0} \leq c$, we have $u(x,t)\leq U(t)$, where $U$ is the solution to the ODE $\begin{cases}
        \frac{d}{dt}U=F(U)\\ U(0)=c
    \end{cases}$.
\end{prop}

\begin{lem}
    Let $g_{t,p}$ be a solution to the NFRF with surface fibers of genus $g\geq 2$, then there exists a constant $C_0$ such that for each $(t,p)$ we have\thinspace:
    $$|K_{g_t}(x)-k| \leq C_0 e^{kt}$$ 
\end{lem}

\begin{proof}
    For each $p\in S$ we can solve the time-dependent Poisson equation 
    $$\Delta_{g_{t,p}} \Phi_p(t) = k-K_{g_{t,p}}$$
    and the solution is well-defined up to a function $f$ depending only on time. This choice can be fixed by requiring that $\int \Phi_p(t) d\text{vol}_{g_{t,p}} = 0$ for every $(t,p)$. We will calculate the derivative $\partial_t \Delta_{g_{t,p}} \Phi_p(t)$ in two different ways (one using the chain rule and another using the equation for $\Phi_p(t)$). The curvature evolution equation under the Normalized Ricci flow is well-known\thinspace:
    $$\partial_t (k-K_{g_{t,p}}(x))=\Delta_{g_{t,p}}(k-K_{g_{t,p}}(x))+K_{g_{t,p}}(x)(k-K_{g_{t,p}}(x))$$
    and we can re-insert the equation for $\Phi_p(t)$ to get
    $$\Delta_{g_{t,p}}\Delta_{g_{t,p}}\Phi_p(t)-(\Delta_{g_{t,p}}\Phi_p(t))^2+k\Delta_{g_{t,p}}\Phi_p(t)$$
    Now for the chain rule calculation, derivative of the Laplacian is also well-known and gives
    $$\partial_t \Delta_{g_{t,p}} \Phi_p(t)=(K_{g_{t,p}}(x)-k)\Delta_{g_{t,p}}\Phi_p(t)+\Delta_{g_{t,p}}\partial_t\Phi_p(t)$$
    which once again using the equation for $\Phi_p(t)$ can be transformed into
    $$-(\Delta_{g_{t,p}}\Phi_p(t))^2+\Delta_{g_{t,p}}\partial_t\Phi_p(t)$$
    Now comparing the two expressions we deduce that 
    $$\Delta_{g_{t,p}}(\partial_t\Phi_p(t)-\Delta_{g_{t,p}}\Phi_p(t)-k\Phi_p(t))=0.$$
    As harmonic functions on a closed manifold are constant, we have 
    $$\partial_t\Phi_p(t)-\Delta_{g_{t,p}}\Phi_p(t)-k\Phi_p(t)=c(t)$$
    where $c(t)$ is a function depending only on time. We will make use of this calculation in a moment. \par 
    We can use the function $\Phi_p(t)$ to construct another auxiliary function 
    $$H_p(t,x):= K_{g_{t,p}}(x)-k + 2|\nabla^{g_{t,p}}\Phi_p(t,x)|^2$$
    which we want to use the maximum principle to. We start by calculating its time derivative\thinspace:
    $$
        \partial_t H_p= \partial_t(K_{g_{t,p}}(x)-k + 2|\nabla^{g_{t,p}}\Phi_p(t,x)|_{g_{t,p}}^2)
    $$
    For the first term, we once again use the curvature evolution equation:
    $$\partial_t (K_{g_{t,p}}(x)-k)=\Delta_{g_{t,p}}(K_{g_{t,p}}(x)-k)+K_{g_{t,p}}(x)(K_{g_{t,p}}(x)-k)$$
    By definition of $\Phi_p$ we can transform the last term into 
    $$K_{g_{t,p}}(x)(K_{g_{t,p}}(x)-k) = (\Delta_{g_{t,p}} \Phi_p(t))^2+k(K_{g_{t,p}}(x)-k)$$
    Now we turn to calculating the time derivative of the $|\nabla^{g_{t,p}}\Phi_p(t,x)|_{g_{t,p}}^2$ term, which we can think of as $g^{-1}_{t,p}(d\Phi_p, d\Phi_p)$\thinspace:
    \begin{multline*}
        \partial_t |\nabla^{g_{t,p}}\Phi_p(t,x)|_{g_{t,p}}^2=-(k-K_{g_{t,p}}(x))|\nabla^{g_{t,p}}\Phi_p(t,x)|_{g_{t,p}}^2+ \\+2g_{t,p}(\nabla^{g_{t,p}}\Phi_p(t,x), \nabla^{g_{t,p}} \partial_t\Phi_p(t,x))
    \end{multline*}
    Using the expression for $\partial_t \Phi_p$ from before, and taking into account that $\nabla^{g_{t,p}} c(t)$ vanishes, we can transform the expression into
    $$(k+K_{g_{t,p}}(x))|\nabla^{g_{t,p}}\Phi_p(t,x)|_{g_{t,p}}^2 +2g_{t,p}(\nabla^{g_{t,p}}\Phi_p(t,x), \nabla^{g_{t,p}} \Delta_{g_{t,p}}\Phi_p(t,x))$$
    Finally, using Bochner's identity to exchange the covariant derivative with the Laplacian
    \begin{multline*}
        2g_{t,p}(\nabla^{g_{t,p}}\Phi_p(t,x), \nabla^{g_{t,p}} \Delta_{g_{t,p}}\Phi_p(t,x))=\Delta_{g_{t,p}}|\nabla^{g_{t,p}}\Phi_p(t,x)|_{g_{t,p}}^2-\\-2|\nabla^{2,g_{t,p}}\Phi_p(t,x)|_{g_{t,p}}^2 -K_{g_{t,p}}|\nabla^{g_{t,p}}\Phi_p(t,x)|_{g_{t,p}}^2
    \end{multline*}
    we obtain 
    $$k|\nabla^{g_{t,p}}\Phi_p(t,x)|_{g_{t,p}}^2+\Delta_{g_{t,p}}|\nabla^{g_{t,p}}\Phi_p(t,x)|_{g_{t,p}}^2-2|\nabla^{2,g_{t,p}}\Phi_p(t,x)|_{g_{t,p}}^2$$
    Combining both calculations yields 
    \begin{multline*}
        \partial_t H_p=\Delta_{g_{t,p}}(K_{g_{t,p}}(x)-k) + (\Delta_{g_{t,p}} \Phi_p(t))^2+k(K_{g_{t,p}}(x)-k) +\\+k|\nabla^{g_{t,p}}\Phi_p(t,x)|_{g_{t,p}}^2+\Delta_{g_{t,p}}|\nabla^{g_{t,p}}\Phi_p(t,x)|_{g_{t,p}}^2-2|\nabla^{2,g_{t,p}}\Phi_p(t,x)|_{g_{t,p}}^2 = \\= \Delta_{g_{t,p}}H_p+kH_p+(\Delta_{g_{t,p}} \Phi_p(t))^2-2|\nabla^{2,g_{t,p}}\Phi_p(t,x)|_{g_{t,p}}^2
    \end{multline*}
    Since $\Delta_{t,p}$ is the trace of the Hessian $\nabla^{2,g_{t,p}}$, we can estimate the last two terms\thinspace:
    $$(\Delta_{g_{t,p}} \Phi_p(t))^2-2|\nabla^{2,g_{t,p}}\Phi_p(t,x)|_{g_{t,p}}^2 \leq 0$$
    which finally yields the inequality
    $$\partial_t H_p \leq \Delta_{g_{t,p}}H_p+kH_p$$
    This leads us to a direct application of the maximum principle (proposition \ref{maxprinc}) with $V=0$ and $F(U)=kU$\thinspace:
    $$K_{g_{t,p}}-k\leq H_p(t) \leq (\underset{x\in F_p}{\max}H_p(0,x)) e^{kt}$$
    Remark that we have not claimed to solve the Poisson equation smoothly in the $p$-variable, so we can't conclude by taking the maximum of $\underset{x\in F_p}{\max}H_p(0,x)$ over $p\in S$ as we do not know that $H_p$ is continuous in $p$. One could probably delve into the theory of 1-parameter Poisson equations on surfaces to prove such result, but we have taken a more straightforward route to find a universal bound. To do so, we proceed by concatenating a series of inequalities\thinspace:
    \begin{enumerate}
        \item by continuity of $K_{g_{t,p}}-k$, the global maximum among all fibers $F_p$ at $t=0$ exists, we will denote it by $k_0$;
        \item the Schauder estimate for $\Delta^{g_{0,p}}$ implies that there is a constant $A_1$ such that 
        $$\underset{x\in F_p}{\max} |\nabla^{g_{0,p}} \Phi_p(x)|^2 \leq A_1 (\underset{x\in F_p}{\max}|\Phi_p(x)|)^2; $$
        \item the Sobolev embedding $H^2 \to L^\infty$ implies the existence of a constant $A_2$ such that
        $$\underset{x\in F_p}{\max} |\Phi_p(x)| \leq A_2 (||\Phi_p||_2 + ||\Delta^{g_{0,p}}\Phi_p||_2)$$
        recall that $\Delta^{g_{0,p}}\Phi_p = K_{g_{t,p}}-k$ and thus $||\Delta^{g_{0,p}}\Phi_p||_2$ can be bounded among all fibers $F$ using continuity;
        \item using the definition of Laplacian first eigenvalue for mean-value zero functions we can deduce 
        $$||\Phi_p||^2_2 \leq \frac{1}{\lambda_1(g_F(0))} ||\nabla^{g_{0,p}} \Phi_p||_2^2.$$
        Along with the standard integration by parts trick and  H\"older's inequality
        $$||\nabla^{g_{0,p}} \Phi_p||_2^2 \leq ||\Phi_p||_2 \cdot ||\Delta^{g_{0,p}} \Phi_p||_2,  $$
        we derive the following inequality (after dividing by $|| \Phi_p||_2$)\thinspace: 
        $$||\Phi_p||_2 \leq \frac{1}{\lambda_1(g_{0,p})} ||\Delta^{g_{0,p}} \Phi_p||_2.$$
        where the first eigenvalue of the Laplacian varies continuously with the fibers, and thus has the global non-zero minimum $\Lambda$.
    \end{enumerate}
    Concatenating these inequalities provides us with the universal bound on $H_p(t=0)$ along all fibers, which we will denote by $H_0^{max}$. This, in turn, gives us the universal bound on the curvature\thinspace: 
    $$K_{g_{t,p}}-k \leq H^{max}_0 e^{kt}$$
    which combined with the bound from the previous lemma lets us conclude\thinspace: 
    $$|K_{g_{t,p}}-k| \leq C_0 e^{kt}$$
    for some constant $C_0$ universal for all fibers. 
\end{proof}

\end{document}
